\chapter*{Conclusion} % L'étoile évite la numérotation mais l'affiche

Dans ce rapport, nous avons posé les fondements théoriques et méthodologiques nécessaires à notre étude comparant le traitement du langage par les modeles basé sur l'architecture des transformers aux cerveau humain. 



Dans un premier temps, notre état de l'art a mis en évidence les similitudes de l'architecture des Transformers avec le modèle neurolinguistique MUC (Mémoire, Unification, Contrôle). Nous avons montré que l'apprentissage prédictif force les modèles à compresser l'information sémantique sur des variétés de basse dimension, un mécanisme que nous supposons similaire a la composante d'unification du MUC.



Sur le plan méthodologique, nous avons justifié le choix de l'algorithme TwoNN afin de mesurer cette compression. Bien que cet algorithme puisse présenter des instabilités sur de très petits échantillons, sa robustesse face aux courbures locales en fait l'outil le plus adapté à notre étude couche par couche. 

Nous avons aussi détaillé notre protocole de génération de stimuli, comparant des phrases sémantiques à des séquences "Jabberwocky". Dans le but de pouvoir isoler l'impact du sens sur la géométrie des représentations.



 Nous avons établi une roadmap détaillée de la phase d'implémentation au cours du semestre prochain, de la validation technique des scripts en février jusqu'à l'analyse critique des résultats en avril.



Cette phase d'implémentation consistera à vérifier l'existence du phénomène de « ramping » et à observer la déformation du profil de Hunchback dans le traitement de phrase dénué de sens.

\addcontentsline{toc}{chapter}{Conclusion et Perspectives}